\subsection{秩 1 投射模}

定理 3. 设 A 是一个环，M 是有限生成的 A-模。

\begin{enumerate}
\def\labelenumi{(\roman{enumi})}
\item
  若存在 A-模 N 使 \(\mathbf{M} \otimes_{\mathbf{A}} \mathbf{N}\) 同构于 A，则模 M 是秩为 1 的投射模。
\item
  反之，若 \(\mathbf{M}\) 是秩为 1 的投射模，\(\mathbf{M}^*\) 是 \(\mathbf{M}\) 的对偶模，则典范同态 \(u: \mathbf{M} \otimes_{\mathbf{A}} \mathbf{M}^* \to \mathbf{A}\) （对应于典范双线性型 \((x, x^*) \to \langle x, x^* \rangle\) ，后者定义在 \(\mathbf{M} \times \mathbf{M}^*\) 上；《代数》第 II 章 §2 第 3 节）是双射的。
\item
  需要证明：对 A 的每个极大理想 m，\(A_{m}\) -模 \(M_{m}\) 是秩为 1 的自由模（定理 2 (b)）；我们可以把 A 换成 \(A_{m}\) ，因而可设 A 是局部环（ \(\S\) 2, 第 7 节, 命题 18）。设 k = A/m。同构 v: \(M \otimes_{A} N \to A\) 定义了一个同构
\end{enumerate}

\[
v \otimes 1 _ {k} \colon (\mathbf {M} / \mathfrak {m M}) \otimes_ {k} (\mathbf {N} / \mathfrak {m N}) \rightarrow k
\]

由于在 \(k\) 上，(M/mM) \(\otimes_k\) (N/mN) 的秩是 M/mM 与 N/mN 的秩之积，后两者必然都等于 1，换言之，M/mM 是单生成的。由此推出 M 是单生成的（§ 3, 第 2 节, 命题 4 的推论 2）；另一方面，M 的零化子也零化 M \(\otimes_A\) N，因而为零，这就证明了 M 同构于 A。

\begin{enumerate}
\def\labelenumi{(\roman{enumi})}
\setcounter{enumi}{1}
\tightlist
\item
  只须证明：对 A 的每个极大理想 m，\(u_{m}\) 是同构（ \(\S 3\) , 第 3 节, 定理 1）。由于 M 是有限表示的（第 I 章, \(\S 2\) , 第 8 节, 引理 8），\((\mathbf{M}^{*})_{\mathfrak{m}}\) 典范地等同于对偶模 \((\mathbf{M}_{\mathfrak{m}})^{*}\) （ \(\S 2\) , 第 7 节, 命题 19），并且由于 \(M_{m}\) 像其对偶模 \((\mathbf{M}_{\mathfrak{m}})^{*}\) 一样是秩为 1 的自由模，显然典范同态 \(u_{\mathfrak{m}}: (\mathbf{M}_{\mathfrak{m}}) \otimes_{\mathbf{A}\mathfrak{m}} (\mathbf{M}_{\mathfrak{m}})^{*} \to \mathbf{A}_{\mathfrak{m}}\) 是双射的，证毕。
\end{enumerate}

注 (1). 若 \(M\) 是秩为 1 的投射模，且 \(\mathbf{N}\) 使 \(\mathbf{M} \otimes_{\mathbf{A}} \mathbf{N}\) 同构于 \(\mathbf{A}\) ，则 \(\mathbf{N}\) 同构于 \(\mathbf{M}^*\) ：有一列同构

\[
\mathbf {N} \rightarrow \mathbf {N} \otimes \mathbf {A} \rightarrow \mathbf {N} \otimes \mathbf {M} \otimes \mathbf {M} ^ {*} \rightarrow \mathbf {A} \otimes \mathbf {M} ^ {*} \rightarrow \mathbf {M} ^ {*}.
\]

命题 6. 设 M 与 N 是秩为 1 的投射 A-模。则 M ⊗\_\{A\} N、Hom\_\{A\}(M, N) 以及 M 的对偶 M* 都是秩为 1 的投射模。

这由公式 (2)、(3) 与 (4) 立即得出。

现在注意，每个有限生成 A-模都同构于 \(L = A^{(N)}\) 的一个商模；因此可以谈论有限生成 A-模按同构关系所分的类组成的集合 \(F(A)\) （《集合论》第 I 章 §6 第 9 节）；我们用 \(P(A)\) 表示 \(F(A)\) 中由秩为 1 的投射 A-模的类组成的子集，并用 \(\operatorname{cl}(M)\) 表示秩为 1 的投射 A-模 M 在 \(P(A)\) 中的像。显然，对两个投射 A-模 M、

N（秩为 1），\(\operatorname{cl}(\mathbf{M} \otimes_{\mathbf{A}} \mathbf{N})\) 只依赖于 \(\operatorname{cl}(\mathbf{M})\) 与 \(\operatorname{cl}(\mathbf{N})\) ；我们定义

\[
\operatorname{cl} (\mathbf {M}) + \operatorname{cl} (\mathbf {N}) = \operatorname{cl} (\mathbf {M} \otimes_ {\mathbf {A}} \mathbf {N})\tag{6}
\]

这样就在 \(\pmb{P}(\mathbf{A})\) 上定义了一个内合成律。

命题 7. 秩为 1 的投射 A-模的类所成的集合 P(A)，连同合成律 (6)，是一个交换群。若 M 是秩为 1 的投射 A-模，M* 是它的对偶模，则

\[
\operatorname{cl} (\mathbf {M} ^ {*}) = - \operatorname{cl} (\mathbf {M}) \quad a n d \quad \operatorname{cl} (\mathbf {A}) = 0.\tag{7}
\]

张量积的结合性与交换性表明合成律 (6) 是结合的和交换的；\(A \otimes_{A} M\) 与 M 之间的同构证明 \(\operatorname{cl}(A)\) 是这个律下的单位元，并且依据定理 3，\(\operatorname{cl}(M) + \operatorname{cl}(M^{*}) = \operatorname{cl}(A)\) ，命题得证。

设 B 是交换 A-代数，M 是秩为 1 的投射 A-模；则 \(\mathbf{M}_{(\mathbf{B})} = \mathbf{B} \otimes_{\mathbf{A}} \mathbf{M}\) 是秩为 1 的投射 B-模（第 3 节, 命题 4）。于是存在一个称为典范的映射 \(\phi: P(\mathbf{A}) \to P(\mathbf{B})\) ，使得

\[
\phi (\operatorname{cl} (M)) = \operatorname{cl} (M _ {(B)}).\tag{8}
\]

对两个 A-模 M、N 成立的等式 \(\mathbf{M}_{(\mathbf{B})}\otimes_{\mathbf{B}}\mathbf{N}_{(\mathbf{B})}=(\mathbf{M}\otimes_{\mathbf{A}}\mathbf{N})_{(\mathbf{B})}\) 证明映射 \(\phi\) 是一个交换群同态。

*注 (2). 定理 1 的条件 (e)（它等价于 P 是投射的且有限生成的这一事实）也可以表述为：与 P 相伴(*)的模层 \(\tilde{\mathbf{P}}\) （它是 \(\mathbf{X} = \operatorname{Spec}(\mathbf{A})\) 上的层） 是局部自由的且是有限型的，因而可以解释为 X 上某个向量丛的截面层。反之，X 上的每个向量丛都来自一个有限生成的投射模，而且该模在相差一个唯一同构的意义下被确定；这样，秩为 \(n\) 的投射模就对应于所有纤维都是 \(n\) 维的向量丛。特别地，秩为 1 的向量丛对应于秩为 1 的投射模。若用 \(\mathcal{O}_{\mathbf{x}}\) 表示结构层 \(\tilde{\mathbf{A}}\) ，用 \(\mathcal{O}_{\mathbf{x}}^{*}\) 表示 \(\mathcal{O}_{\mathbf{x}}\) 的可逆元层（它在 X 的开集 U 上的截面是 \(\mathcal{O}_{\mathbf{x}}\) 在 U 上截面环中的可逆元），则由此得出群 \(P(\mathbf{A})\) 同构于第一上同调群 \(\mathrm{H}^1 (\mathrm{X},\mathcal{O}_\mathrm{x}^*)\) 。

